% =====================================================================
%  c20 —— 原文印刷页 165（PDF 第 166 页）
%  含：第 10 节末尾「关于任意子的更多材料」文献列表的续接部分
%  （第 2--4 条与 OTHER REFS 速记清单），以及 Section 11
%  （11.1、11.2、11.3、11.4）的全部条目。
%  本片段无编号公式、无插图。
% =====================================================================

% 下列编号列表是从上一页（印刷 p.164）接续下来的：上一页在
% 「Here are some sources for more material about anyons:」之下已给出第 1 条
% （John Preskill 的讲义），本页从第 2 条起排，故用 start=2 复刻原编号。
\begin{enumerate}[start=2]
  \item 若要找一份相当新的综述，可参见 A. Stern，「Anyons and the quantum Hall effectA pedagogical review」，Annals of Physics 323: 204；arXiv:0711.4697v1。
  \item A. Lerda，\textit{Anyons: Quantum mechanics of particles with fractional statistics}，Lect. Notes Phys. M14 (1992) 1-138。
  \item 其它可参考的文献还包括：A. Khare，\textit{Fractional Statistics and Quantum Theory}；以及 G. Dunne，\textit{Self-Dual Chern-Simons Theories}。
\end{enumerate}

% [Tutor Note: 原文第 2 条里「effectA pedagogical review」处丢了一个句读——该综述（arXiv:0711.4697）的正式标题是「Anyons and the quantum Hall effect: A pedagogical review」，PDF 文本层与页面图像里冒号都被吞掉了（连字 ff 的排版所为）。此处照原样转写为 effectA，不加补；抄录文献时请以 arXiv 标题为准。另：原文该处用的是英式 ”…” 双引号，译文按本工程规范统一为「」。]

\textbf{其它文献（OTHER REFS）：}
\begin{itemize}[label={}]
  \item S. Buron，「A Short Guide to Anyons and Modular Functors」，
  \item 「From categories to anyons: a travelogue」
  \item Preskill
  \item Rowell 与 Wang：[52] BAMS
\end{itemize}

% [Tutor Note: 这块「OTHER REFS」是作者随手记下的补充读物清单，与上面那个编号列表不属于同一层，故单独排成无标号列表。四条都只给了残缺信息：「S. Buron」的拼写疑为笔误，「From categories to anyons: a travelogue」缺作者与出处，「Preskill」即上一页第 1 条已给出的 John Preskill 讲义，「Rowell and Wang: [52] BAMS」中的 [52] 是本文参考文献表的编号、BAMS 指 Bulletin of the American Mathematical Society。译文一律照原样保留，不作补全，以免替作者「猜出」文献。]

\section{进一步推广与应用的综述} % 原文: 11. A Survey Of Further Generalizations And Applications

\begin{keytakeaways}
  \item 本节是作者给出的一张「应用地图」：只列条目，不给推导，每条线索背后都是独立的一支研究文献。
  \item 超引力方向四条：RR 场与自对偶性、反常抵消、M 理论的 Chern-Simons 项，以及带边流形上 M 理论配合 $E_8$ 规范理论时的反常抵消。
  \item 非紧群方向：三维量子引力（Brown-Henneaux 与 Witten 的进路），以及把能级 $k$ 解析延拓、改用复群。
  \item 全息方向：$SL(2,\mathbb{R})$ 与单态表示（singleton）。
  \item 其它方向：格点费米子（David Kaplan）、算术 Chern-Simons（Minhyong Kim）、超群上的 Chern-Simons、与拓扑弦的关系（Marino 的综述 [39]）、CY 3-流形上的全纯 Chern-Simons 对应 B 模型开拓扑弦，以及弦场论。
\end{keytakeaways}

\subsection{在超引力中的应用} % 原文: 11.1 Applications To Supergravity

\begin{enumerate}
  \item RR 场与自对偶性
  \item 反常抵消。
  \item M 理论的 Chern-Simons 项。
  \item 带边流形上 M 理论（配合 $E_8$ 规范理论）的反常抵消。
\end{enumerate}

% [Tutor Note: 第一条的「RR 场」是 Ramond--Ramond 场（原文只写缩写），即 II 型超引力中的那些反对称张量场；「自对偶性」指其场强满足一阶自对偶约束 $*F=\pm F$（例如 IIB 五形式场强），这类约束无法由普通的二阶（度规型）作用量变分得到，必须改用 Chern-Simons 型拓扑项来写，这正是 CS 理论在此处登场的原因。本节四条都只是标题，作者未展开任何一条，故此处不补推导。]

\subsection{非紧群} % 原文: 11.2 Noncompact Groups

三维量子引力：
\begin{itemize}[label={}]
  \item Brown-Henneaux
  \item Witten 处理三维量子引力的诸种进路。
  \item 在 $k$ 上的解析延拓：复群。
\end{itemize}

% [Tutor Note: 「在 $k$ 上的解析延拓」指把 CS 能级 $k$ 当作复参数：对紧群，$k$ 必须是整数（作用量的 $2\pi$ 整性、以及标架反常给出的平移 $k\to k+h^\vee$ 都需要它），而一旦允许复 $k$、乃至把规范群复化，这套整性条件与度规的正定假设都要重新检查——这是非紧群条目下最易被忽略的一句。作者在此只点了线索，细节需查 Witten 关于三维引力的原文。]

\subsection{在全息中的应用} % 原文: 11.3 Applications to Holography

$\mathrm{SL}(2,\mathbb{R})$

单态表示（singletons）

% [Tutor Note: 原文这两行分别印作直立的「SL(2,R)」与「Singletons」，此处把 R 排成 $\mathbb{R}$ 只是记号规范化。术语上 singleton（Dirac 的「单态」表示，即 $SO(2,d)$ 的超轻表示）若直译成「单态」易与 $J=0$ 的那种单态混淆，故保留英文并附中文；下文若再出现同一词，不再重复括注。]

\subsection{许多其它应用与推广} % 原文: 11.4 Many Other Applications And Extensions

\begin{enumerate}
  \item 格点费米子（David Kaplan）
  \item 算术 Chern-Simons（Minhyong Kim）
  \item 超群上的 Chern-Simons
  \item Chern-Simons 与拓扑弦。参见 Marino 的综述 [39]。
  \item CY 3-流形上的全纯 Chern-Simons：B 模型开拓扑弦
  \item 弦场论
\end{enumerate}

% [Tutor Note: 第五条极易读错：「CY 3-folds」是复三维 Calabi-Yau 流形，实维数为 6，因此其上的全纯 Chern-Simons 是一个六维作用量，与前面三维 CS 不是同一个理论，不能把三维的 $k$ 量子化、Wilson 线等结论直接搬来；「B 模型」指拓扑 B 模型，其开弦（也就是拓扑弦）由全纯 CS 描述，这一点正是它被列入「推广」而非「应用」的原因。]
